%%%%%%%%%%%%%%%%%%%%%%%%%%%%%%%%%%%%%%%%%

\documentclass[paper=letter, fontsize=12pt]{article}
\usepackage[english]{babel} % English language/hyphenation
\usepackage{amsmath,amsfonts,amsthm} % Math packages
\usepackage[utf8]{inputenc}
\usepackage{float}
\usepackage{lipsum} % Package to generate dummy text throughout this template
\usepackage{blindtext}
\usepackage{graphicx} 
\usepackage{caption}
\usepackage{xcolor}
\usepackage{enumerate}
\usepackage{eso-pic}
\usepackage{tikz}
\usepackage{subcaption}
\usepackage[hidelinks]{hyperref}
\usepackage[sc]{mathpazo} % Use the Palatino font
\usepackage[T1]{fontenc} % Use 8-bit encoding that has 256 glyphs
\linespread{1.05} % Line spacing - Palatino needs more space between lines
\usepackage{microtype} % Slightly tweak font spacing for aesthetics
\usepackage[hmarginratio=1:1,top=32mm,columnsep=20pt]{geometry} % Document margins
\usepackage{multicol} % Used for the two-column layout of the document
%\usepackage[hang, small,labelfont=bf,up,textfont=it,up]{caption} % Custom captions under/above floats in tables or figures
\usepackage{booktabs} % Horizontal rules in tables
\usepackage{float} % Required for tables and figures in the multi-column environment - they need to be placed in specific locations with the [H] (e.g. \begin{table}[H])
\usepackage{hyperref} % For hyperlinks in the PDF
\usepackage{lettrine} % The lettrine is the first enlarged letter at the beginning of the text
\usepackage{paralist} % Used for the compactitem environment which makes bullet points with less space between them
\usepackage{abstract} % Allows abstract customization
\renewcommand{\abstractnamefont}{\normalfont\bfseries} % Set the "Abstract" text to bold
\renewcommand{\abstracttextfont}{\normalfont\small\itshape} % Set the abstract itself to small italic text
\usepackage{titlesec} % Allows customization of titles

\usepackage{parskip}
\renewcommand\thesection{\Roman{section}} % Roman numerals for the sections
\renewcommand\thesubsection{\Roman{subsection}} % Roman numerals for subsections

\definecolor{myRed}{HTML}{660000}

\graphicspath{ {images/} }
\newcommand\BackgroundPic{%
	\begin{tikzpicture}[remember picture, overlay, opacity=0.08]%
	\node [rotate=0,scale=0.8,text opacity=0.79] at (current page.center) {\includegraphics{./logo_10}};
	\end{tikzpicture}
}

\titleformat{\section}[block]{\Large\scshape\color{myRed}\bf}{\thesection.}{1em}{} % Change the look of the section titles
\titleformat{\subsection}[block]{\large}{\thesubsection.}{1em}{} % Change the look of the section titles
\newcommand{\horrule}[1]{\rule{\linewidth}{#1}} % Create horizontal rule command with 1 argument of height
\usepackage{fancyhdr} % Headers and footers
\pagestyle{fancy} % All pages have headers and footers
\fancyhead{} % Blank out the default header
\fancyfoot{} % Blank out the default footer
\author{}
\date{}
\definecolor{myBlue}{HTML}{000066}
\hypersetup{
	colorlinks,
	linkcolor={myRed!50!black},
	citecolor={myRed!50!black},
	urlcolor={myRed!80!black}
}

\fancyhead[C]{University of Central Flordia $\bullet$ Spring 2017 $\bullet$ IDS6938 - Simulation Techniques } % Custom header text

\fancyfoot[RO,LE]{\thepage} % Custom footer text
%----------------------------------------------------------------------------------------
%       TITLE SECTION
%----------------------------------------------------------------------------------------
\title{\vspace{-15mm}\fontsize{24pt}{10pt}\selectfont\textbf{Simulation Techniques (IDS6938) }\\ \fontsize{15pt}{10pt}\selectfont\textbf{Modeling \& Simulation} | Graduate Program\\University of Central Florida\\\fontsize{11pt}{10pt}\selectfont Spring 2017- T,R 6:00pm-7:20pm} % Article title


%----------------------------------------------------------------------------------------
\begin{document}
\AddToShipoutPicture{\BackgroundPic} 
\maketitle % Insert title
\thispagestyle{fancy} % All pages have headers and footers


\vspace{-8mm}
\section{Instructor Contact:}
\vspace{-4mm}
\hspace{25pt} Instructor: Dr. Joseph T. Kider Jr. \\
\hspace*{25pt} Office: Partnership II: Office 320 \\
\hspace*{25pt} Office Hours: Tuesday, Thursday 5pm-6pm\\
\hspace*{25pt} Email: \href{mailto:jkider@ist.ucf.edu}{jkider@ist.ucf.edu}

\section{Overview:}
This course provides a broad survey of three different categories of simulation techniques: continuous simulation, discrete event simulation, and agent-based simulation. Material in this course will provide an overview of the foundations as well as specific example problems and simulation solutions for each category of simulation. Students will be expected to understand the basics for a variety of different types of simulations, as well as to work with hands-on simulation tools to implement solutions various problems.\\

\noindent \textbf{Prerequisites}: DIG5876 - Quantitative Aspects of Modeling \& Simulation or\textit{ Instructor Consent}\\
\textbf{Recommended}: A grasp of \textit{Calculus} and \textit{Programming}

\section{Learner Objectives:}
The	main objectives	of this course include showing the students how continuous, discrete event, and agent-based simulation techniques appear in real life and physical phenomena. This requires understanding the fundamental mathematics and understanding how to apply these methods to numerical simulation. By the end of the course, the students should be able to make mathematical simulation models involving differential equations for problems encountered in engineering, social and physical sciences, and to solve them by using one or a combination of the methods mentioned above. Therefore we will teach students the basics of analytic, geometric and numerical methods, for studying differential equations for simulation. We will also teach students the fundamentals of programming the numerical methods for each type of simulation technique.

This is a very ambitious course and will push students. Be prepared for hard gratifying work.

This course will be very practical and prepare students for problems they will encounter.



\section{Course Text Material}
\textcolor{myRed}{\textbf{Required text:}} NONE. This is a ``unique" course. There is no one text that covers all material completely. The instructor will provide resources on CANVAS to supplement learning from a variety of texts and resources.

\textcolor{myRed}{\textbf{Recommended text (sources heavily used in class):}}
\begin{itemize}\setlength\itemsep{-0.6em}
	\item Simulation, Modeling \& Analysis (Law)
	\item Discovering Modern C++: An Intensive Course for Scientists, Engineers, and Programmers (C++ In-Depth)
	\item Partial Differential Equations: Second Edition (Evans)
	\item Partial Differential Equations: An Introduction (Strauss)
	\item Partial Differential Equations for Scientists and Engineers
	\item Div, Grad, Curl, and All That: An Informal Text on Vector Calculus (Schey).
	\item The Big Book of Simulation Modeling. Multimethod Modeling with AnyLogic (Borshchev)
	\item"AnyLogic 7 in Three Days (Grigoryev)
	\item Discrete-Event System Simulation (Bank et. al)
\end{itemize}
\textit{(Again- I will provide needed chapters of these works on CANVAS.)}


\section{Course Topics}

The course will include the following topics. (\textit{Note: Subject to change since this is the first semester teaching the course. The instructor reserves the right to slow or speedup the material as deemed appropriate. This is a rough guide of topics.})

\textbf{Continuous simulation}
\begin{itemize}\setlength\itemsep{-0.6em}
\item Overview of differential equations: ordinary differential equations, partial differential equations, and system dynamics \& stability concepts.
\item Decomposition \& discretization: fixed interval timing, dynamic interval timing, mesh and particle-based methods.
\item Overview of numerical solvers: forward \& reverse Euler, Runga-Kutta, finite volume vs. finite element methods
\item Numerical Error analysis: taylor series, analytical, etc
\end{itemize}

\textbf{Discrete event simulation}
\begin{itemize}\setlength\itemsep{-0.6em}
	\item Queuing models: single-server models, queuing networks
	\item Basics of DES: event based modeling, time-advance \& optimistic timing models, distributed simulations
	\item DES analytical methods: input and output analysis, monte carolo sampling
	\item Advanced Sampling techniques: importance sampling, error analysis
\end{itemize}

\textbf{Agent-based modeling and simulation}
\begin{itemize}\setlength\itemsep{-0.6em}
	\item Game Theory: normal vs extensive form games, pure vs. mixed, min-max and regret, Nash equilibria \& rationality
	\item Communication and networks: signalling and aggregation, voting methods, auction systems
	\item Simple ABMs: nature-based flocking, herding, and following, cellular automta
\end{itemize}

\section{Grading \& Course Requirements:}
The grade distribution is as follows: final exam ($25\%$), projects ($75\%$).
\\\\
\noindent \textcolor{myRed}{\textbf{Exam:}} There will be a final exam in the course, worth $25\%$ of the total grade. The final exam submission will be due at the time assigned by UCF.  The final exam is open book and will be available for 1 week so you may take it at any point between the time it opens and it is due. \\\\
\textcolor{myRed}{\textbf{Projects:}} There will be three projects, one for each category of simulation. These projects will be individual work. Students will provide implemented simulation solutions intended to demonstrate a synthesis and understanding of using simulation to solve realistic problems. The projects will be posted online via UCF’s CANVAS system, and the solutions will be submitted
via CANVAS, as well by 11:59pm on Wednesday. You will be provided roughly 2 weeks per assignment. You will get 3 \textit{late days} to use at your discretion during the semester. Homework turned in within 24 hours after the deadline are worth half their original points, after that it \textbf{will not be graded}.
\\\\
\textcolor{myRed}{\textbf{Letter Grades:}}
Letter grades will  be determined at the end of the semester by finding ``natural" divisions between clusters of numeric grades based roughly on the table shown below. 



\begin{table}[h!]
	\centering
	\label{my-label}
	\begin{tabular}{|l|l|p{95mm}|}\hline
		\textbf{Letter Grade} & \textbf{Points} & \textbf{Description} \\ \hline
		\textbf{A} & 90-100 & Outstanding work! Much more than would normally be expected for a student at this level. Very strong performance. I learned something valuable from you. \\ \hline
		\textbf{B} & 80-89 & Good work! Clearly above average. More than would normally be expected for a student at this level. Above average work and performance. Very few minor flaws.\\ \hline
		\textbf{C} & 70-79 & OK. Directions were followed. You did what would normally be expected for the	assignment for a student at this level. A few minor flaws. \\ \hline
		\textbf{D} & 60-69 & You can do much better. You have performed at a level that is below what one would expect for a student at this level. A number of minor flaws and/or one major flaw. \\ \hline
		\textit{E} & 59-60 & \textit{(This is a made up grade to make sure you are still reading but makes you think why we skip E?) }\\ \hline
		\textbf{F} & 59 or below & What happened? Well below average. More than one major flaw and/or an inordinate	number of minor flaws. \\ \hline
  		
	\end{tabular}
\end{table}

\textit{The $+$ and $-$ modifiers for grades will be determined at my discretion, largely based on class participation.}

\textcolor{myRed}{\textbf{Grade postings:}} All grades will be posted on CANVAS.

\textcolor{myRed}{\textbf{Make up exams:}} As the projects / exam are online, the students may take them any time up until their final due date. If you will not be
able to take the test on the day it is due, you must plan to take it prior to that time. In the case of an emergency, contact the instructor immediately to arrange for alternative plans

\textcolor{myRed}{\textbf{Attendance Policy:}} According to the university, this class is "mixed mode". Here is the official description:
 \textbf{M – Mixed Mode/Reduced Seat Time} courses include both required classroom attendance and online instruction. Classes have substantial activity conducted over the web, which substitutes for some classroom meetings.
 
 This will be a ``\textbf{Flipped Classroom}" class. Students are expected to show up prepared. This explicitly means watching the online module material \textbf{before} class. During class you are expected to work in teams to solve problems, code examples, and ask questions. \textbf{Attendance it is extremely important to success in this class and should be considered mandatory}.


\textcolor{myRed}{\textbf{Academic Engagement:}} All instructors/faculty are required to document students' academic activity at the beginning of each course. In order to
document that you began this course, please complete the following academic activity by the end of the first week of classes or as soon as possible after adding the course. Failure to do so may result in a delay in the disbursement of your
financial aid. To document your academic engagement, you will complete a survey quiz by the end of the first week of class. This survey quiz must be completed in order to unlock the remainder of the curriculum. This quiz has a practical purpose to help guide your learning throughout the semester.

\section{Academic Honesty:}
Plagiarism and cheating of any kind on an examination, quiz, or assignment will result at least in an "F" for that assignment (and may, depending on the severity of the case, lead to an "F" for the entire course) and may be subject to appropriate referral to the \href{http://osc.sdes.ucf.edu/faq#students}{Office of Student Conduct} for further action. See the \href{http://goldenrule.sdes.ucf.edu/}{UCF Golden Rule} for further information. I will assume for this course that you will adhere to the academic creed of this University and will maintain the highest standards of academic integrity. In other words, don't cheat by giving answers to others or taking them from anyone else. I will also adhere to the highest standards of academic integrity, so please do not ask me to change (or expect me to change) your grade illegitimately or to bend or break rules for one person that will not apply to everyone. Do not cheat. NOTE: As graders, we reserve the right to compare content submitted by students to alternative content
(either students’ or otherwise) by any means necessary (automated or non-automated).


Use of cell phone and talking during class, which results in the disruption of other students, are not allowed. Please silence cell phones before class. Laptops are highly encouraged (\textit{required one per team}). Students should follow the instructor lessons and refrain from inappropriate material. Use of websites like \textit{Stackoverflow} and \textit{programming subreddits} are highly encouraged.

\section{Accessibility Statement:}
The University of Central Florida is committed to providing reasonable accommodations for all persons with disabilities. This syllabus is available in alternate formats upon request. Students with disabilities who need accommodations in this course must contact the professor at the beginning of the semester to discuss needed accommodations. No accommodations will be provided until the student has met with the professor to request accommodations. Students who need accommodations must be registered with Student Accessibility Services, Ferrell Commons, 7F, Room 185, phone (407) 823-2371, TTY/TDD only phone (407) 823-2116, before requesting accommodations from the professor.

\section{Changes/Announcements:}
Though I try to plan for everything, life will surprise us and we may have to make adjustments. In the event that anything in this syllabus changes (e.g. classroom moves, changes in due dates, contact information), I will use a broad announcement so that all students are informed immediately.\textbf{ It is critically important that you set your CANVAS announcements so that you receive all notifications and be sure to check the class site at least twice a week.}

\section{Copyright:}
This course may contain copyright protected materials such as audio or video clips, images, text materials, etc. These items are being used with regard to the Fair Use doctrine in order to enhance the learning environment. Please do not copy, duplicate, download or distribute these items. The use of these materials is strictly reserved for this online classroom environment and your use only. All copyright materials are credited to the copyright holder.

\section{Third-Party Software and FERPA}
During this course you might have the opportunity to use public online services and/or software applications sometimes called third-party software and new media. Do not post or provide any private information about yourself or your classmates. Where appropriate you may use a pseudonym or nickname. Some assignments posted publicly may require personal reflection/comments, but the assignments will not require you to disclose any personally identity-sensitive information. If you have any concerns about this, please contact me.


%----------------------------------------------------------------------------------------
%\end{multicols}
\end{document}
